\documentclass[11pt]{article}

\ifdefined\pdftrailerid
  \pdftrailerid{}
\fi

\usepackage[a4paper,margin=30mm]{geometry}
\usepackage[T1]{fontenc}
\usepackage{lmodern}
\usepackage{amsmath,amssymb,amsthm}
\usepackage{booktabs}
\usepackage{enumitem}
\usepackage{microtype}
\usepackage{xcolor}
\usepackage{xurl}
\usepackage[colorlinks=true,linkcolor=blue!55!black,citecolor=blue!55!black,
  urlcolor=blue!55!black]{hyperref}
\hypersetup{
  pdftitle={AI-generated Research Dossier - Odd-Cycle-Free Latin Squares: Products, Congruence Obstructions, and an Order-8 Census},
  pdfauthor={},
  pdfsubject={A compact structural and computational study of odd-cycle-free Latin squares},
  pdfkeywords={Latin squares, quasigroups, Latin trades, direct products, congruences}
}

\title{Odd-Cycle-Free Latin Squares:\\
  Products, Congruence Obstructions, and an Order-$8$ Census}
\author{}
\date{\small Preserved technical manuscript; read the preceding provenance statement}

\newtheorem{theorem}{Theorem}[section]
\newtheorem{lemma}[theorem]{Lemma}
\newtheorem{corollary}[theorem]{Corollary}
\newtheorem{proposition}[theorem]{Proposition}
\newtheorem{computationaltheorem}[theorem]{Computational Theorem}
\theoremstyle{definition}
\newtheorem{definition}[theorem]{Definition}
\newtheorem{conjecture}[theorem]{Conjecture}
\theoremstyle{remark}
\newtheorem{remark}[theorem]{Remark}

\newcommand{\pat}{\operatorname{pat}}
\newcommand{\rowv}{\mathrm{row}}
\newcommand{\colv}{\mathrm{col}}
\newcommand{\symv}{\mathrm{sym}}
\newcommand{\FFF}{\mathsf{FFF}}
\newcommand{\TTT}{\mathsf{TTT}}

\newcommand{\DossierVariant}{Compact reading version}
\begin{document}
\pagenumbering{roman}
\thispagestyle{empty}
\noindent{\large AI-generated mathematical research dossier\par}
\vspace{1.1cm}
\noindent{\LARGE\bfseries Odd-Cycle-Free Latin Squares\par}
\vspace{0.45cm}
\noindent{\Large \DossierVariant\par}
\vspace{0.7cm}
\noindent{\large Public materials for critical examination\par}
\vspace{0.8cm}

\noindent\textbf{Origin.}
The mathematical development, proof text, software and scientific writing
were produced through OpenAI ChatGPT/Codex workflows, using the cited human
literature, census data and software. The human project role was initiation
and provision of time and computing resources, not subject-matter expertise
or a claimed personal scientific contribution.

\medskip
\noindent\textbf{Review and responsibility.}
No accountable scholarly author or independent expert human review of the
complete work has been established. AI systems are disclosed as tools, not
listed as authors. This is a research dossier for examination, not a
peer-reviewed publication or a conventional author-approved paper.

\medskip
\noindent\textbf{Evidence boundary.}
The order-$10$ exclusion is classified in the research record as a written
reduction plus exact computation, not a solver-free proof or a
proof-assistant certificate. The global power-of-two conjecture remains
open; order $12$ is the next undecided even non-power-of-two case in this
project. Automated checks are not substitutes for expert review or a
literature-based novelty assessment.

\vfill
\noindent\textbf{Public dossier, 19 September 2026.}
The compact and full versions overlap; they are not independent studies.
No DOI, journal acceptance, institutional endorsement or new licence grant
is asserted. Read the scope statement before using the technical text.
\clearpage

\section*{Reading, Evidence and Distribution Scope}
\subsection*{How to read the technical text}
First-person expressions such as ``we prove'' belong to the generated
exposition. The theorem labels describe the evidence claimed by the research
record, not an external referee's endorsement. References to independent
implementations describe specific computational checks, not independent
human expert review of the entire work. Attribution to the cited human
authors, data providers and software developers remains intact.

The mathematical sections and bibliography are preserved from the technical
record. Title-page metadata, this disclosure and availability wording are
editorial changes only. The compact version omits substantial parts of the
long record, including its full order-$10$ argument.

\subsection*{Public export, not the complete working archive}
This release includes the two dossiers, their editable TeX sources and a
curated source-and-result collection. It does not include personal
correspondence, development history, software environments or large graph
and checkpoint archives. External census files must be obtained from their
cited providers. The public README and reproducibility guide define exactly
which checks can be run from this export.

References in the technical text to historical packages describe the
underlying research record. They do not assert that every historical
artifact or dependency is included in this curated export. Files containing
machine-local paths have a documented privacy projection; public file hashes
therefore need not equal historical whole-file hashes. Numerical results
are not recomputed or upgraded by that projection. The public manifest and
projection report distinguish these cases.

\subsection*{Computational limits}
Executed computations provide stronger evidence than generated prose alone,
but audit success establishes only the checks actually performed. Census
completeness and the primitive-group classification remain cited external
dependencies. Differently sliced runs of the same search binary are not
independent implementations. No new exhaustive order-$10$ search is claimed
by publishing these materials. A complete public, persistent data deposit
for the multi-gigabyte master graphs and checkpoints is still outstanding.

\subsection*{Rights and intended use}
Rights remain reserved pending an explicit licensing decision, to the extent
rights exist. Public hosting does not establish copyright in generated
content, grant third-party rights or imply an open-source licence.
The purpose is to make the material inspectable and invite specific
corrections. No claim of novelty, scholarly qualification or personal
mathematical authorship is inferred from the hosting account.
\clearpage
\pagenumbering{arabic}



\begin{center}
{\LARGE Odd-Cycle-Free Latin Squares:\par}
\vspace{0.35em}
{\LARGE Products, Congruence Obstructions, and an Order-$8$ Census\par}
\vspace{1.25em}
{\small Preserved technical manuscript\par}
\vspace{0.7em}

{\small Read the preceding provenance statement\par}
\end{center}
\vspace{0.75em}

\begin{abstract}
For each pair of rows, columns, or symbols of a Latin square, the corresponding
two-line permutation is fixed-point-free.  We call a view failing when one of
these permutations has a nontrivial odd cycle, and call the square FFF when
none of its three views fails.  Such cycles support Latin trades.  We prove
the exact three-parity sign table: odd row- and column-view trades reverse the
conventional row--column sign, while odd symbol-view trades preserve it and
reverse the two parastrophic companion signs.

We prove that odd-order squares are TTT, characterize FFF group tables, and
give a basis-family closure criterion for group isotopy.  Failure patterns are
monotone from Latin subsquares and are preserved exactly through congruences
with two-element blocks; odd block size forces TTT.  Consequently every FFF
quasigroup of order $2p$, for odd prime $p$, is congruence-simple in every
isotope.  A twisted binary extension gives an explicit non-group-isotopic FFF
loop of order $8$.  Nonuniversal binary line dependencies are in bijection
with half-order subsquares and occur in complementary four-packs.  Hence an
FFF square of order $2$ modulo $4$, greater than $2$, has maximal binary
line-incidence rank.  Specializing known $3$-net character theory,
nonuniversal prime-field line dependencies are balanced cyclic quotient
homotopies.  Hence at order $2p$, for odd prime $p$, one F view forces a
saturated integer line-incidence lattice.

For direct products we prove
\[
  \pat(L\times K)=\pat(L)\mathbin{\mathrm{OR}}\pat(K)
\]
componentwise, and show that $L\times K$ is group-isotopic exactly when both
factors are.  An exhaustive reproducible scan of the $283657$ order-$8$ main
classes finds $230$ FFF classes, of which $225$ are not group-isotopic; all
eight patterns occur.  Exact generation also shows that no order-$6$ FFF
square exists.  Hence every order $2^m$, $m\geq3$, realizes all eight patterns
and contains a non-group-isotopic FFF square.  These results motivate, but do
not prove, the conjecture that nontrivial FFF squares exist exactly in
power-of-two orders.
\end{abstract}

\noindent\textbf{Keywords.}
Latin square; quasigroup congruence; Latin subsquare; Latin trade; direct
product; computational enumeration.

\medskip
\noindent\textbf{2020 Mathematics Subject Classification.}
Primary 05B15; Secondary 20N05, 05C70.

\section*{Evidence convention}

Unlabelled lemmas, propositions, and theorems below are proved mathematically
from the stated definitions and cited results.  Statements labelled
\emph{Computational Theorem} are exhaustive finite computations with frozen
inputs, executable code, commands, hashes, and independent validation reports.
Any corollary that explicitly invokes a Computational Theorem inherits that
computational dependency.  These results are not presented as formal
proof-assistant certificates.  The complete artifact map is given in
Section~\ref{sec:reproducibility}.

The trade-parity convention is explicit: the conventional row--column sign
and its two parastrophic companions are distinguished.  No symbol-view trade
is claimed to reverse the conventional sign merely by parastrophy.

\section{Introduction}

An odd cycle in the permutation induced by two lines of a Latin square
supports a simple trade: swap the two lines on that cycle.  The move preserves
the Latin property.  In the row and column views it reverses the conventional
row--column Latin sign; in the symbol view it preserves that sign and reverses
the two parastrophically transported signs.  This local operation is therefore
a natural primitive in parity-based approaches related to the Alon--Tarsi
setting, but its three sign conventions must be kept distinct.  B\'erczi's
experimental program studies such local corrections as ingredients for
global constructions~\cite{Berczi2026}; it does not supply the theorems proved
here.

We isolate the underlying existence question.  For each of the row, column,
and symbol views, ask whether some pair of lines induces a permutation with a
nontrivial odd cycle.  The three answers form a pattern in $\{F,T\}^3$.  A
square is FFF when all three answers are negative.  Equivalently, every union
of two link factors in the tripartite representation has component lengths
divisible by four.

Our first contribution is structural.  Every odd-order Latin square is TTT,
whereas the Cayley table of a finite group is FFF exactly when the group is a
$2$-group.  A basis-family closure criterion detects group isotopy.  We then
show that failure patterns are monotone from Latin subsquares and analyze
quasigroup congruences: odd block size forces TTT, while two-element blocks
preserve the quotient pattern exactly.  Thus a hypothetical FFF quasigroup of
order $2p$, for odd prime $p$, must be congruence-simple in every isotope.
This excludes all congruence-imprimitive constructions at order $10$ but does
not by itself decide the simple case.

The remaining binary line-incidence defect also has an exact subsquare
meaning.  Nonuniversal vectors in the left binary line kernel are in
bijection with half-order subsquares, and the quotient classes group them
into complementary four-packs.  This yields a closed counting formula and
forces maximal binary line rank for every FFF order congruent to $2$ modulo
$4$ and greater than $2$.  At order $10$ this is necessary, not sufficient.

More generally, Moorhouse's $3$-net character formula specializes as follows:
a nonuniversal line dependency over a prime field is a
balanced cyclic quotient homotopy, and after isotopy its fibres form a
quasigroup congruence.  Consequently every Smith torsion prime divides the
order.  For order $2p$, with $p$ odd prime, the congruence obstruction shows
that even one F view forces a saturated integer line-incidence lattice.

The group construction is not the whole phenomenon.  A twisted binary
extension of $C_4$ gives an explicit non-group-isotopic FFF loop of order $8$.
Moreover, direct products satisfy the exact componentwise formula
\[
  \pat(L\times K)=\pat(L)\mathbin{\mathrm{OR}}\pat(K),
\]
and a product is group-isotopic exactly when both factors are.  These theorems
turn one order-$8$ example into infinite power-of-two families.

Our computational contribution is exhaustive.  McKay, Meynert, and Myrvold
enumerated the main classes of small Latin squares~\cite{McKayMeynertMyrvold2007};
the associated ANU file contains one representative of each of the $283657$
order-$8$ main classes~\cite{McKayLatinData}.  Our scan finds exactly $230$
FFF main classes.  Only $5$ are group-isotopic and $225$ are not, and all
eight ordered failure patterns occur.  Independently, fresh exact-cover
generation of all $9408$ reduced order-$6$ squares finds no FFF square and
matches the official ANU set exactly.

Combining theory and computation proves that every order $2^m$, $m\geq3$,
realizes all eight patterns and contains a non-group-isotopic FFF square.  We
do not claim that the product construction is injective on isotopy or main
classes.  We also do not prove the Alon--Tarsi conjecture.  The accompanying
complete research report proves by exhaustive master searches that order $10$
has no FFF square; that computation is outside this compact paper's theorem
scope.  The resulting open conjecture is narrower: nontrivial FFF Latin
squares may exist exactly in power-of-two orders.

The longer companion research report in the repository develops additional
cycle-parity, subset-action, primitive-group, and flag-surface machinery
and proves nonexistence at order $10$.  None of that machinery is needed for
the results in this compact paper.

\section{Definitions and invariance}

Let $X$ be a set of size $n$.  A \emph{Latin square} on $X$ is a map
$L\colon X\times X\to X$ such that every row and every column is a bijection.
For $a,c,u\in X$, define the row, column, and symbol-line bijections
\[
 r_a(c)=L(a,c),\qquad
 q_c(a)=L(a,c),\qquad
 s_u(c)=\text{the unique }a\text{ such that }L(a,c)=u.
\]
For two lines in the corresponding view, put
\[
 \rho_{a,b}=r_b^{-1}r_a,\qquad
 \kappa_{c,d}=q_d^{-1}q_c,\qquad
 \sigma_{u,v}=s_v^{-1}s_u.
\]
All three permutations act on an $n$-element complementary coordinate set.

\begin{lemma}\label{lem:fixed-point-free}
If the two lines are distinct, the corresponding induced permutation is
fixed-point-free.
\end{lemma}

\begin{proof}
A fixed point of $\rho_{a,b}$ would give $L(a,c)=L(b,c)$ in one column,
contrary to the Latin property.  The column case is identical.  A fixed point
of $\sigma_{u,v}$ would say that the same cell contains the two distinct
symbols $u$ and $v$.
\end{proof}

\begin{definition}
A view of $L$ \emph{fails} if one of its induced two-line permutations has an
odd cycle of length greater than one.  Write
\[
 \pat(L)=(\rowv(L),\colv(L),\symv(L))\in\{F,T\}^3,
\]
where $T$ means that the view fails.  A square is \emph{FFF} if
$\pat(L)=\FFF$.
\end{definition}

The terminology ``fails'' is chosen from the perspective of the local-trade
program: a $T$ view supplies an odd-cycle trade, while an $F$ view does not.

\begin{definition}
Fix two rows $a\ne b$ and a union $C$ of cycles of $\rho_{a,b}$.  The
associated \emph{two-line trade} swaps the entries in rows $a,b$ in precisely
the columns in $C$.  Column- and symbol-view trades are defined by permuting
the coordinate roles.
\end{definition}

\begin{proposition}\label{prop:trade}
A two-line trade produces another Latin square.  Define the three line-sign
products
\[
 R(L)=\prod_a\operatorname{sgn}(r_a),\qquad
 C(L)=\prod_c\operatorname{sgn}(q_c),\qquad
 S(L)=\prod_u\operatorname{sgn}(s_u),
\]
and put
\[
 \epsilon_{RC}=RC,\qquad \epsilon_{RS}=RS,\qquad \epsilon_{CS}=CS.
\]
The first is the conventional row--column Latin sign, while the other two
are its parastrophic companions.  If a trade has support size $k$ and
$\delta=(-1)^k$, then its sign ratios are
\[
\begin{array}{c|ccc|ccc}
\text{trade view} & R'/R&C'/C&S'/S&
 \epsilon'_{RC}/\epsilon_{RC}&\epsilon'_{RS}/\epsilon_{RS}&
 \epsilon'_{CS}/\epsilon_{CS}\\
\hline
\rowv    &1&\delta&\delta&\delta&\delta&1\\
\colv    &\delta&1&\delta&\delta&1&\delta\\
\symv    &\delta&\delta&1&1&\delta&\delta
\end{array}
\]
Consequently an odd row- or column-view trade reverses the conventional
$\epsilon_{RC}$, whereas an odd symbol-view trade preserves $\epsilon_{RC}$
and reverses $\epsilon_{RS}$ and $\epsilon_{CS}$.
\end{proposition}

\begin{proof}
Write $\pi=\rho_{a,b}$.  Since $r_a(c)=r_b(\pi(c))$ and $C$ is
$\pi$-invariant, the two rows contain the same symbol set on $C$.  Swapping
those entries therefore preserves both row symbol sets, while each affected
column merely swaps two entries.  Thus the result is Latin.

Extend $\pi|_C$ by the identity outside $C$.  The two changed row
permutations are obtained by precomposition with $\pi$ and $\pi^{-1}$, so
their sign changes cancel and $R'/R=1$.  Each of the $|C|$ affected column
permutations is composed with a transposition, so $C'/C=(-1)^{|C|}$.  The
same support contains the same $|C|$ symbols in the two rows; each associated
symbol-position permutation swaps those two rows, giving
$S'/S=(-1)^{|C|}$.  This proves the first table row.  Exchanging rows and
columns proves the second.

For a symbol-view trade, each support column swaps the two selected symbols,
and the cycle relation pairs the affected cells in exactly the same number of
rows.  Hence $R'/R=C'/C=(-1)^{|C|}$.  The two changed symbol-position
permutations differ from their predecessors by mutually inverse
permutations on the support, so their sign changes cancel and $S'/S=1$.
Multiplying the appropriate line-sign ratios proves the last three columns.
\end{proof}

\begin{proposition}[Tripartite link formulation]
\label{prop:tripartite-links}
Regard the cells of $L$ as the triangle decomposition
\[
 \bigl\{\{a,c,L(a,c)\}:a,c\in X\bigr\}
\]
of the complete tripartite graph on the row, column, and symbol copies of
$X$.  The links of the vertices in each part form a 1-factorization between
the other two parts.  An induced two-line cycle of length $\ell$ corresponds
bijectively to an alternating cycle of length $2\ell$ in the union of the two
associated link factors.  Hence $L$ is FFF if and only if every two-factor
union in all three link factorizations has only cycle lengths divisible by
four.

Furthermore, the total number of $4$-cycles among the two-factor unions is
the same in all three views and equals the number of intercalates of $L$.
\end{proposition}

\begin{proof}
For a row $a$, its incident triangles give the perfect matching
$M_a=\{\{c,r_a(c)\}:c\in X\}$ between columns and symbols.  For distinct
rows $a,b$, alternately following an edge of $M_a$ and an edge of $M_b$
sends a column $c$ to $r_b^{-1}r_a(c)=\rho_{a,b}(c)$.  Thus an $\ell$-cycle
of $\rho_{a,b}$ is exactly one alternating $2\ell$-cycle of
$M_a\cup M_b$.  Cyclically permuting the three coordinate classes proves the
column and symbol statements and the FFF equivalence.

A $4$-cycle in $M_a\cup M_b$ is exactly a $2$-cycle of $\rho_{a,b}$, hence a
$2$-by-$2$ Latin subsquare on two rows, two columns, and two symbols.  This is
an intercalate.  The same intercalate is counted once by its row pair, once by
its column pair, and once by its symbol pair, proving the three equalities.
\end{proof}

\begin{lemma}[Global three-view sign identity]
\label{lem:global-three-view-sign}
For every Latin square of order $n$,
\[
 R(L)C(L)S(L)=(-1)^{\binom n2}.
\]
\end{lemma}

\begin{proof}
Give each coordinate copy of $X$ the same fixed order, and let
$\tau(x,y)=(y,x)$ on $X^2$.  The $\binom n2$ off-diagonal pairs are the
transpositions of $\tau$, so
$\operatorname{sgn}(\tau)=(-1)^{\binom n2}$.

The row-entry map $A(r,c)=(r,L(r,c))$ is block diagonal by rows and has sign
$R(L)$.  The map $B(r,c)=(c,L(r,c))$ is a coordinate transposition followed
by the block-diagonal family of column permutations, so its sign is
$(-1)^{\binom n2}C(L)$.  Likewise, $D(c,s)=(s,s_s(c))$ has sign
$(-1)^{\binom n2}S(L)$.  Since $A=\tau\circ D\circ B$, taking signs gives
\[
 R(L)=(-1)^{3\binom n2}S(L)C(L)
      =(-1)^{\binom n2}S(L)C(L),
\]
which is equivalent to the claimed identity.
\end{proof}

For later comparison, let $C_{v,k}$ denote the total number of $k$-cycles
among all induced line-pair permutations in view $v$.

\begin{proposition}[Aggregate cycle-profile baseline]
\label{prop:aggregate-profile-baseline}
Let $n$ be even and let $p$ be prime.  If
\[
 \sum_{k=2}^n a_kx_k+b=0\pmod p
\]
for every cycle-count vector $(x_2,\ldots,x_n)$ of a fixed-point-free
permutation of degree $n$, then, separately for each view $v$,
\[
 \sum_{k=2}^n a_k C_{v,k}+\binom n2 b=0\pmod p.
\]
Moreover,
\[
 C_{\rowv,2}=C_{\colv,2}=C_{\symv,2}
\]
and
\[
 \sum_{v\in\{\rowv,\colv,\symv\}}\ \sum_{\substack{2\leq k\leq n\\k\text{ even}}}
 C_{v,k}=\frac{n(n-1)}2\pmod2.
\]
\end{proposition}

\begin{proof}
The aggregate profile in one view is the sum of the
$\binom n2$ fixed-point-free line-pair profiles, proving the first identity.
The second identity is the intercalate bijection from
Proposition~\ref{prop:tripartite-links}.  For the last identity, the sign of
a fixed-point-free permutation of even degree is $(-1)^e$, where $e$ is its
number of even cycles.  Multiplying relative-permutation signs over all line
pairs in a view gives the product of its line signs.
Lemma~\ref{lem:global-three-view-sign} gives the product over the three views as
$(-1)^{n(n-1)/2}$, which proves the congruence.
\end{proof}

\begin{proposition}[Line-pair sign cut]\label{prop:line-sign-cut}
Fix one view and write its line bijections as $\lambda_i$.  For
$P_{ij}=\lambda_j^{-1}\lambda_i$, the pairs for which $P_{ij}$ is odd form
a cut of the complete graph on the lines.  If $n$ is even and the view is
F, then the pairs for which $P_{ij}$ has an odd number of cycles form the
same cut.  Consequently, for even $n>2$, no view can have every induced
line-pair permutation equal to one $n$-cycle.
\end{proposition}

\begin{proof}
Put $\varepsilon_i=\operatorname{sgn}(\lambda_i)$.  Since sign is a
homomorphism,
\[
 \operatorname{sgn}(P_{ij})=\varepsilon_i\varepsilon_j.
\]
Thus $P_{ij}$ is odd exactly when the endpoint signs differ.  These pairs
are the edges of the cut between positive and negative lines.

In an F view of even degree, every cycle of $P_{ij}$ has even length and
therefore negative sign.  Hence
$\operatorname{sgn}(P_{ij})=(-1)^{c(P_{ij})}$, where $c(P_{ij})$ is the
number of cycles.  Finally, if all relative permutations were $n$-cycles,
every edge would be negative.  A triangle of three lines would then cross
the sign cut three times, whereas every cut meets a cycle evenly.
\end{proof}

An \emph{isotopy} independently permutes rows, columns, and symbols.  A
\emph{parastrophy} permutes the three coordinate roles.  Their combined
orbits are the \emph{main classes} (also called species).

\begin{proposition}\label{prop:invariance}
Under isotopy, each of the three view bits is invariant.  Under parastrophy,
the three bits are permuted.  Consequently, being FFF and having at least one
failing view are main-class invariants.
\end{proposition}

\begin{proof}
Let an isotopy act by permutations $\alpha,\beta,\gamma$ of rows, columns,
and symbols.  The transformed row maps satisfy
$r'_{\alpha(a)}=\gamma r_a\beta^{-1}$, and hence
\[
 (r'_{\alpha(b)})^{-1}r'_{\alpha(a)}
   =\beta(r_b^{-1}r_a)\beta^{-1}.
\]
Thus row cycle lengths are unchanged.  The column calculation conjugates by
$\alpha$, and the symbol calculation conjugates by $\beta$.  Parastrophy
only permutes which coordinate pair is called row, column, or symbol.
\end{proof}

\section{Elementary structural results}

\subsection{Odd orders and group tables}

\begin{proposition}\label{prop:odd-order}
Every Latin square of odd order $n>1$ has pattern $\TTT$.
\end{proposition}

\begin{proof}
By Lemma~\ref{lem:fixed-point-free}, every induced permutation between two
distinct lines is fixed-point-free on an odd set.  If all of its nontrivial
cycles were even, their lengths would have even sum.  Hence it has an odd
cycle of length greater than one.  This argument applies separately to any
pair of distinct rows, columns, and symbols.
\end{proof}

Let $G$ be a finite group and let $L_G(x,y)=xy$ be its Cayley table.

\begin{proposition}\label{prop:group-tables}
The row permutation induced by rows $a,b$ of $L_G$ is left multiplication by
$b^{-1}a$.  The column permutation induced by columns $c,d$ is right
multiplication by $cd^{-1}$.  The symbol permutation induced by symbols
$u,v$ is right multiplication by $u^{-1}v$.  Each cycle length is therefore
the order of the corresponding group element.
\end{proposition}

\begin{proof}
The row identity is
$r_b^{-1}r_a(x)=b^{-1}ax$.  Likewise
$q_d^{-1}q_c(x)=xcd^{-1}$.  Finally, $s_u(c)=uc^{-1}$ and
$s_v^{-1}(a)=a^{-1}v$, whence
$s_v^{-1}s_u(c)=cu^{-1}v$.  Left or right multiplication by $g$ has orbits
equal to cosets of $\langle g\rangle$, all of length $\operatorname{ord}(g)$.
\end{proof}

\begin{corollary}\label{cor:group-fff}
The Cayley table of a finite group $G$ is FFF if and only if $G$ is a
$2$-group.
\end{corollary}

\begin{proof}
If $G$ is a $2$-group, every nonidentity element has even order, so
Proposition~\ref{prop:group-tables} gives no nontrivial odd cycle.  Conversely,
if $|G|$ has an odd prime divisor, Cauchy's theorem supplies a nonidentity
element of odd order; the corresponding induced permutations give a failing
view.
\end{proof}

\subsection{A closure criterion for group isotopy}

Fix a base row $b$ of a Latin square $L$ and define
\[
 H_L(b)=\{r_b^{-1}r_x:x\in X\}\subseteq\operatorname{Sym}(X).
\]

\begin{proposition}[Basis-family closure criterion]\label{prop:closure}
A Latin square $L$ is isotopic to a group table if and only if $H_L(b)$ is
closed under composition.  This condition is independent of the chosen base
row.
\end{proposition}

\begin{proof}
Suppose first that $H=H_L(b)$ is closed.  It contains the identity, has
$|X|$ distinct elements, and is a finite submonoid of a symmetric group;
hence it is a group.  For each $z\in X$, the values
$(r_b^{-1}r_x)(z)$, as $x$ varies, are all distinct, since a column of $L$
contains every symbol once.  Thus $H$ acts regularly on $X$.

Fix $z_0\in X$.  Relabel rows by $x\mapsto h_x=r_b^{-1}r_x$.  Use the regular
action to identify every column $z$ with the unique $h_z\in H$ satisfying
$h_z(z_0)=z$, and identify every relabeled symbol $r_b^{-1}(u)$ in the same
way.  The entry in row $h_x$ and column $h_z$ is represented by the unique
element sending $z_0$ to $h_xh_z(z_0)$, namely $h_xh_z$.  The resulting
isotope is therefore the Cayley table of $H$.

Conversely, for a group table the family at base row $b$ is the full left
regular representation $\{\lambda_{b^{-1}x}:x\in G\}$ and is closed.
Isotopy conjugates the family and relabels its members, preserving closure.
The same argument for any base row proves base independence.
\end{proof}

The same criterion can equivalently be written in the column or symbol view,
by parastrophy.

\section{Subsquare and congruence obstructions}

The FFF property is hereditary for Latin subsquares in a stronger,
componentwise sense.

\begin{theorem}[Subsquare monotonicity]\label{thm:subsquare-monotonicity}
If $M$ is a Latin subsquare of $L$, then
\[
  \pat(M)\leq \pat(L)
\]
componentwise, with $F<T$.
\end{theorem}

\begin{proof}
Let the rows, columns, and symbols of $M$ be $R_0,C_0,S_0$.  For
$a,b\in R_0$, the ambient row permutation $r_b^{-1}r_a$ preserves $C_0$:
if $c\in C_0$, then $L(a,c)\in S_0$, and row $b$ contains this symbol at a
unique column of $C_0$.  Its restriction is precisely the induced row
permutation of $M$.  Thus every odd row cycle of $M$ is also an odd row
cycle of $L$.  Interchanging rows and columns proves the column statement.

For $u,v\in S_0$ and $c\in C_0$, the row $s_u(c)$ lies in $R_0$.  In that
row, the unique occurrence of $v$ lies in a column $d\in C_0$.  Hence
$s_v^{-1}s_u$ preserves $C_0$ and restricts to the symbol permutation of
$M$.  This proves the symbol statement.
\end{proof}

\begin{corollary}\label{cor:no-odd-subsquare}
Every Latin subsquare of an FFF square is FFF.  In particular, an FFF square
contains no subsquare of odd order greater than one.
\end{corollary}

\begin{proof}
The first assertion is Theorem~\ref{thm:subsquare-monotonicity} with
$\pat(L)=\FFF$.  The second follows from Proposition~\ref{prop:odd-order}.
\end{proof}

Together with Computational Theorem~\ref{cthm:small-orders}, the same implication
excludes order-$6$ subsquares; that additional exclusion inherits the
computational status of the order-$6$ theorem.

There is an exact linear-algebraic description of all half-order
subsquares.  Let $M_L$ be the $3n$-by-$n^2$ binary incidence matrix whose
rows are the row, column, and symbol lines and whose columns are the occupied
cells of $L$.  Let
\[
 U=\{(\alpha\mathbf1_R,\beta\mathbf1_C,\gamma\mathbf1_S):
       \alpha+\beta+\gamma=0\}
 \subseteq\ker M_L^{\mathsf T}
\]
and put
\[
 \delta_2(L)=(3n-2)-\operatorname{rank}_{\mathbb F_2}M_L.
\]

\begin{theorem}[Binary line kernel and half-order subsquares]
\label{thm:binary-kernel-subsquares}
Individual embedded half-order Latin subsquares of $L$ are in bijection with
\[
 \ker M_L^{\mathsf T}\setminus U.
\]
Equivalently, the nonzero classes in
$\ker M_L^{\mathsf T}/U$ parametrize their complementary four-packs.  Thus
their number is exactly
\begin{equation}
 4\bigl(2^{\delta_2(L)}-1\bigr).\label{eq:half-order-count}
\end{equation}
If $n=2m$, where $m>1$ is odd, then any one F-view forces
\begin{equation}
 \operatorname{rank}_{\mathbb F_2}M_L=3n-2.\label{eq:fff-binary-line-rank}
\end{equation}
In particular, \eqref{eq:fff-binary-line-rank} holds for every FFF square of order congruent to $2$
modulo $4$ and greater than $2$.
\end{theorem}

\begin{proof}
A left-kernel vector is a triple $(a,b,c)$ satisfying
\begin{equation}
 a_r+b_j+c_{L(r,j)}=0.\label{eq:binary-line-dependency}
\end{equation}
If one coordinate, say $a$, is constant, then fixing a
column and varying its rows shows that $c$ is constant; fixing a row then
makes $b$ constant.  Hence every vector outside $U$ has all three
coordinates nonconstant.

Fix a row $r$.  Since that row permutes the symbols,
\eqref{eq:binary-line-dependency} identifies the
multiset of $c$-values with that of $a_r+b$.  Choosing rows with $a_r=0$
and $a_r=1$ gives
\[
 \operatorname{wt}(c)=\operatorname{wt}(b)
 =n-\operatorname{wt}(b).
\]
Thus $n=2m$ and $b,c$ each have weight $m$.  The analogous fixed-column
argument gives $\operatorname{wt}(a)=m$.  For their binary fibres
$A_x,B_y,C_z$, equation~\eqref{eq:binary-line-dependency} now gives
\[
 L(A_x,B_y)\subseteq C_{x+y}\qquad(x,y\in\mathbb F_2).
\]
Every set has size $m$; Latinness makes each of the four displayed blocks an
order-$m$ Latin subsquare.

Adding an element of $U$ complements zero or two of the three partitions and
selects a different block in the same four-pack.  Conversely, if
$R_0\times C_0$ is a half-order subsquare with symbol set $S_0$, row and
column counting forces the complementary pattern
\[
 R_0\times C_0,\ R_1\times C_1\longrightarrow S_0,
 \qquad R_0\times C_1,\ R_1\times C_0\longrightarrow S_1.
\]
The three complement indicators satisfy
\eqref{eq:binary-line-dependency}, and the chosen block recovers
them uniquely.  This proves both bijections.  Since $U$ has dimension two
and the quotient has dimension $\delta_2(L)$, counting its nonzero classes
proves the counting formula.

Finally suppose that $m>1$ is odd and $\delta_2(L)>0$.  The preceding
construction gives an order-$m$ subsquare.  Two distinct rows inside it
induce a fixed-point-free permutation preserving its $m$ columns.  Such a
permutation on an odd set has a nontrivial odd cycle, contradicting an F row
view.  The two parastrophes give the same conclusion for either other view,
which proves the rank formula.  The exception $m=1$ is real: the order-$2$ FFF square has
four singleton subsquares and $\delta_2=1$.
\end{proof}

The complete finite audit independently constructs both sides of the
bijection on $9778$ tables and obtains exact set equality in every case.  In
particular, the $230$ order-$8$ FFF representatives have
\[
 \delta_2=0,1,2,3\quad\text{with multiplicities}\quad63,156,10,1,
\]
which explains their half-order-subsquare distribution
$0^{63}4^{156}12^{10}28^1$.  A hypothetical order-$10$ FFF square must have
binary line rank $28$ and no order-$5$ subsquare.  This condition is not
sufficient: all $135$ tracked order-$10$ diagnostic tables already satisfy
it, none is FFF, and that corpus is not exhaustive.

The same kernel equation has a uniform interpretation over every prime
field, now in terms of quotient homotopies rather than individual
subsquares.  Moorhouse's character formula for loop-coordinatized
$3$-nets gives the general identity
$\operatorname{rank}_{p}M_L=3n-2-\dim\operatorname{Hom}(G,\mathbb F_p)$
and identifies its defect with elementary-abelian quotient
characters~\cite[Theorem~4.2]{Moorhouse1991BruckNets}.  We record a direct
three-sorted proof in our notation because the individual cyclic quotient is
the input needed for the FFF corollary below.
The author's publication page flags an erratum for the original article; the
reference entry therefore points to the author-hosted article copy and the
separate erratum listing.  The direct proof given here, rather than an unexamined appeal
to the original pagination, is the logical basis for the specialization used
below.

\begin{theorem}[Prime-field line quotients]
\label{thm:prime-field-line-quotients}
Let $p$ be prime and let
$(a,b,c)\in\ker_{\mathbb F_p}M_L^{\mathsf T}$ be nonconstant modulo the
two-dimensional constant subspace.  Then $p\mid n$, each of $a,b,c$ takes
every value exactly $n/p$ times, and their fibres satisfy
\[
 L(R_x,C_y)\subseteq S_{-x-y}\qquad(x,y\in\mathbb F_p).
\]
After a suitable isotopy these three fibre partitions become one quasigroup
congruence with $p$ blocks of size $n/p$ and quotient isotopic to $C_p$.

Over characteristic zero, $M_L$ has rank $3n-2$.  Every prime dividing a
nonzero Smith invariant factor of the integer matrix $M_L$ divides $n$.
\end{theorem}

\begin{proof}
If one of $a,b,c$ is constant, fixing a line in either other sort and using
Latinness makes the remaining two constant.  Let $\mu_b(t)$ and $\mu_c(t)$
be value multiplicities.  Since every row permutes the symbols,
equation~\eqref{eq:binary-line-dependency}, now read in $\mathbb F_p$, gives
\[
 \mu_c(t)=\mu_b(-a_r-t).
\]
Comparing rows with $a_r-a_{r'}\ne0$ makes $\mu_b$ invariant under a nonzero
translation.  That translation generates the additive group of
$\mathbb F_p$, so $\mu_b$ is uniform.  Thus $p\mid n$ and both $b,c$ are
balanced; the column argument balances $a$.  The displayed fibre rule is the
kernel equation itself.

Choose bijections from one partitioned $n$-set to the row, column, and symbol
sets, matching equally labelled fibres.  Transporting $L$ by these three
bijections gives an isotope whose block operation is $x*y=-x-y$ on
$\mathbb F_p$.  It is Latin and isotopic to addition, so block equivalence is
compatible with multiplication and both divisions and is a quasigroup
congruence.

In characteristic zero a finite nonempty multiset cannot be invariant under
a nonzero translation, because its iterated translates would have infinite
support.  Hence only the two constant dependencies remain.  Finally, if a
prime divides a nonzero Smith invariant factor, reduction modulo that prime
lowers the rank; the first part then forces the prime to divide $n$.
\end{proof}

An eight-worker audit checks this theorem over
$\mathbb F_2,\mathbb F_3,\mathbb F_5,\mathbb F_7,\mathbb F_{11}$ on $9778$
tables.  It verifies all $1160$ nonuniversal quotient classes and a separate
implementation reconstructs $48890$ modular ranks, with no errors.

We next use congruences in the quasigroup algebra with multiplication and the
two division operations.  All congruence blocks have equal size and form a
Latin quotient.

\begin{theorem}[Congruence obstructions]\label{thm:congruence-obstructions}
Let $L$ be a finite Latin square.
\begin{enumerate}[label=(\roman*)]
\item If $L$ has a congruence with odd block size greater than one, then
  $\pat(L)=\TTT$.
\item If the congruence blocks have size two and the quotient is $Q$, then
  $\pat(L)=\pat(Q)$ componentwise.
\end{enumerate}
\end{theorem}

\begin{proof}
For (i), choose distinct rows $x,x'$ in one congruence block $A$ and any
column block $B$.  The two left translations map $B$ bijectively onto the
same product block.  Thus $L_{x'}^{-1}L_x$ preserves $B$ and is
fixed-point-free there.  Since $|B|$ is odd, its restriction has a
nontrivial odd cycle.  Applying the same argument to the two relevant
parastrophes proves the column and symbol bits.

For (ii), label every two-element block by $(i,0),(i,1)$, where $i\in Q$.
The unique form of a binary Latin square gives a function
$f\colon Q^2\to\mathbb F_2$ such that
\[
  (i,a)*(j,b)=(i*j,a+b+f(i,j)).
\]
Rows in the same quotient block induce only transpositions inside the column
fibres.  For rows in different quotient blocks, the lift above a quotient
cycle of length $\ell$ is either two cycles of length $\ell$ or one cycle of
length $2\ell$, according to its binary voltage.  Thus an ambient odd cycle
can occur only above an odd quotient cycle.

Conversely, suppose quotient rows $i,i'$ induce an odd cycle $C$ of length
$\ell>1$.
Above a quotient column $j\in C$, the lifted row permutation comparing
$(i,a)$ with $(i',a')$ has the form
\[
 (j,b)\longmapsto
 \bigl(\pi(j),b+a+a'+f(i,j)+f(i',\pi(j))\bigr).
\]
Its voltage around $C$ is $a+a'+S$, because $\ell$ is odd, where $S$ is
independent of $a,a'$.  Choose $a+a'=S$.  The voltage is then zero, so the
lift of $C$ consists of two cycles of the same odd length $\ell$.  Hence a
quotient row witness lifts to an ambient row witness.  This proves equality
of the row bits.  The same complete argument in the parastrophes proves
equality of the column and symbol bits.
\end{proof}

\begin{corollary}[Order-$2p$ simplicity]\label{cor:2p-simple}
Let $p$ be an odd prime.  Every quasigroup of order $2p$ with at least one
F view is congruence-simple, and every isotope of it is congruence-simple.
\end{corollary}

\begin{proof}
A nontrivial congruence has block-size/quotient-order pair $(p,2)$ or
$(2,p)$.  The first is excluded by
Theorem~\ref{thm:congruence-obstructions}(i).  In the second, the quotient
has pattern TTT by Proposition~\ref{prop:odd-order}, so (ii) excludes it.
Both congruence cases force TTT, contradicting even one F view.
Finally, the full view pattern is isotopy invariant by
Proposition~\ref{prop:invariance}, and the same argument applies to every isotope.
\end{proof}

\begin{corollary}[Order-$2p$ line-incidence saturation]
\label{cor:2p-line-saturation}
Let $p>1$ be an odd prime and let $L$ have order $2p$.  If even one view of
$L$ is F, then $M_L$ has rank $6p-2$ over every field, all its nonzero Smith
invariant factors are $1$, and
\[
 \mathbb Z^{6p}/\operatorname{im}_{\mathbb Z}M_L\cong\mathbb Z^2.
\]
In particular, this holds for every hypothetical order-$10$ FFF square.
\end{corollary}

\begin{proof}
By Theorem~\ref{thm:prime-field-line-quotients}, a torsion prime must divide
$2p$.  A characteristic-two defect yields, after isotopy, a congruence with
two blocks of odd size $p$, so
Theorem~\ref{thm:congruence-obstructions}(i) gives pattern TTT.  A
characteristic-$p$ defect yields blocks of size two; part (ii) identifies the
pattern with that of the odd-order quotient, again TTT by
Proposition~\ref{prop:odd-order}.  Isotopy invariance contradicts the assumed
F view in either case.  Thus no Smith prime remains, proving saturation and
the displayed cokernel.
\end{proof}

Thus a hypothetical order-$10$ FFF square must be isotopically
congruence-simple.  This removes every congruence-imprimitive quasigroup
construction, but it does not rule out simple quasigroups and therefore does
not decide order $10$.

The binary equality iterates through congruence series.

\begin{corollary}[Prime-factor congruence series]
\label{cor:prime-congruence-series}
Suppose that a quasigroup admits a chain of congruence quotients ending in the
one-element quasigroup, with every congruence block size prime.  It is FFF if
all those primes are $2$, and it has pattern TTT if at least one is odd.
\end{corollary}

\begin{proof}
A binary step preserves the entire pattern by
Theorem~\ref{thm:congruence-obstructions}(ii).  At an odd step,
Theorem~\ref{thm:congruence-obstructions}(i) gives TTT independently of the
lower quotient.  Iterating proves both statements.
\end{proof}

This proves the power-of-two conjecture for this congruence-decomposable
class, but leaves congruence-simple quasigroups outside its scope.

The same theorem also gives a census-free nongroup example.

\begin{theorem}[Explicit nongroup FFF loop]
\label{thm:explicit-nongroup-loop}
On $\mathbb Z_4\times\mathbb F_2$, define
\[
 (i,a)*(j,b)=
 \bigl(i+j,\,a+b+\mathbf 1_{(i,j)=(1,1)}\bigr),
\]
with coordinates taken modulo $4$ and $2$, respectively.  This is an FFF
loop of order $8$ that is not group-isotopic.
\end{theorem}

\begin{proof}
The identity is $(0,0)$, and the first-coordinate fibres form a congruence
with quotient $C_4$.  The quotient is FFF by
Corollary~\ref{cor:group-fff}, so binary pattern equality makes the loop
FFF.  It is not associative because
\[
 ((1,0)*(1,0))*(2,0)=(0,1),\qquad
 (1,0)*((1,0)*(2,0))=(0,0).
\]
At the identity base row, the basis family in
Proposition~\ref{prop:closure} is the family of left translations.
Closure would force $L_xL_y=L_{x*y}$ and hence associativity.  The displayed
failure is therefore a closure failure, so the loop is not group-isotopic.
\end{proof}

\section{Direct products}

Let $L$ be a Latin square on $X$ and $K$ a Latin square on $Y$.  Their direct
product $P=L\times K$ is the Latin square on $X\times Y$ defined by
\[
 P((x,y),(x',y'))=(L(x,x'),K(y,y')).
\]

\begin{lemma}[Product-cycle formula]\label{lem:product-cycle}
Let $\pi$ and $\tau$ be permutations.  If $x$ lies on a $\pi$-cycle of
length $m$ and $y$ lies on a $\tau$-cycle of length $n$, then $(x,y)$ lies on
a cycle of length $\operatorname{lcm}(m,n)$ under $\pi\times\tau$.
\end{lemma}

\begin{proof}
The pair returns after $t>0$ steps exactly when both $m\mid t$ and $n\mid t$.
The least such $t$ is $\operatorname{lcm}(m,n)$.
\end{proof}

\begin{theorem}[Pattern product theorem]\label{thm:pattern-product}
For arbitrary Latin squares $L$ and $K$,
\[
 \pat(L\times K)=\pat(L)\mathbin{\mathrm{OR}}\pat(K)
\]
componentwise.
\end{theorem}

\begin{proof}
Rows of $P$ are indexed by $(a,\alpha)$, and their row maps are
$r^P_{(a,\alpha)}=r^L_a\times r^K_\alpha$.  Hence the induced permutation
between $(a,\alpha)$ and $(b,\beta)$ is
\[
 \rho^P_{(a,\alpha),(b,\beta)}
   =\rho^L_{a,b}\times\rho^K_{\alpha,\beta}.
\]
If one coordinate of the two product rows is equal, the corresponding factor
is the identity.  The identical calculation for column maps gives
$\kappa^P=\kappa^L\times\kappa^K$.

For the symbol view,
\[
 s^P_{(u,\eta)}(c,d)=(s^L_u(c),s^K_\eta(d)),
\]
and therefore $\sigma^P=\sigma^L\times\sigma^K$.  Thus the same argument
applies in all three views.

If a factor has a nontrivial odd cycle in one view, keep the other factor line
fixed.  The identity contributes cycle length $1$, so the product contains
the same odd cycle.  Conversely, suppose a product induced permutation has
an odd cycle of length $\ell>1$.  By Lemma~\ref{lem:product-cycle},
$\ell=\operatorname{lcm}(m,n)$ for factor cycle lengths $m,n$.  Both are odd,
and at least one is greater than one.  That factor already has a nontrivial
odd cycle in the same view.  This proves the Boolean OR formula in each view.
\end{proof}

\begin{corollary}\label{cor:fff-product}
$L\times K$ is FFF if and only if both $L$ and $K$ are FFF.  In particular,
product with a finite $2$-group table preserves the full pattern of the other
factor.
\end{corollary}

\begin{proof}
The first assertion is the all-$F$ case of
Theorem~\ref{thm:pattern-product}.  For the second, combine the same theorem
with Corollary~\ref{cor:group-fff}.
\end{proof}

\begin{theorem}[Group-isotopy product theorem]\label{thm:group-product}
$L\times K$ is group-isotopic if and only if both $L$ and $K$ are
group-isotopic.
\end{theorem}

\begin{proof}
At a product base row $(b,c)$, the basis family factors as
\[
 H_{L\times K}(b,c)=H_L(b)\times H_K(c).
\]
If both factor families are closed, their Cartesian product is closed.
Conversely, both factor families contain the identity.  If the product family
is closed, then $(h_1,\mathrm{id})(h_2,\mathrm{id})$ lies in it for all
$h_1,h_2\in H_L(b)$, proving closure of $H_L(b)$; using
$(\mathrm{id},k_i)$ proves closure of $H_K(c)$.  Apply
Proposition~\ref{prop:closure}.
\end{proof}

\begin{corollary}\label{cor:inflation}
Let $L$ be a non-group-isotopic FFF Latin square of order $n$.  For every
$m\geq0$, the product $L\times C_2^m$ is a non-group-isotopic FFF Latin
square of order $n2^m$.
\end{corollary}

\begin{proof}
The table of $C_2^m$ is FFF by Corollary~\ref{cor:group-fff}; apply
Corollary~\ref{cor:fff-product} and
Theorem~\ref{thm:group-product}.
\end{proof}

\section{Computational results}

The following statements are exhaustive finite computations.  Their frozen
inputs, programs, commands, outputs, and hashes are specified in
Section~\ref{sec:reproducibility}.

\begin{computationaltheorem}[Small orders]\label{cthm:small-orders}
The unique reduced Latin square of order $2$ is FFF, and all four reduced
Latin squares of order $4$ are FFF.  None of the $9408$ reduced Latin squares
of order $6$ is FFF.  More precisely, their patterns are
\[
 6600\,\TTT,\qquad 936\,\mathsf{TFF},\qquad
 936\,\mathsf{FTF},\qquad 936\,\mathsf{FFT}.
\]
\end{computationaltheorem}

The exact-cover generator loads no cached square list.  Every generated square
was checked to be Latin and reduced, all $9408$ table strings were distinct,
and the generated set agrees exactly with the ANU reduced-order-$6$ data
file~\cite{McKayLatinData}.  Proposition~\ref{prop:invariance} therefore
extends nonexistence from reduced squares to all Latin squares of order $6$.

\begin{computationaltheorem}[Order-$8$ pattern census]\label{cthm:order8}
Across the $283657$ order-$8$ main-class representatives, the pattern
distribution is
\begin{center}
\begin{tabular}{cr@{\qquad}cr@{\qquad}cr@{\qquad}cr}
\toprule
Pattern & Count & Pattern & Count & Pattern & Count & Pattern & Count\\
\midrule
$\FFF$ & 230 & $\mathsf{FFT}$ & 81 & $\mathsf{FTF}$ & 40 &
$\mathsf{FTT}$ & 315\\
$\mathsf{TFF}$ & 133 & $\mathsf{TFT}$ & 209 & $\mathsf{TTF}$ & 232 &
$\TTT$ & 282417\\
\bottomrule
\end{tabular}
\end{center}
Among the $230$ FFF main classes, exactly $5$ are group-isotopic and $225$
are not group-isotopic.
\end{computationaltheorem}

\begin{computationaltheorem}[Subsquare census of the FFF classes]
Among the $230$ FFF representatives, the number of intercalates ranges from
$2$ to $112$.  Each endpoint is attained by exactly one representative: local
source line $10776$ (ANU record $10775$) is the unique minimum class and is not
group-isotopic, while local line $1$ (ANU record $0$) is the unique maximum
class and is group-isotopic.  Local physical lines are numbered from one;
ANU records are numbered from zero.  No representative
contains an order-$3$ subsquare.  The distribution of the number of order-$4$
subsquares is
\[
  0^{63},\qquad 4^{156},\qquad 12^{10},\qquad 28^1,
\]
where exponents count representatives.
\end{computationaltheorem}

The subsquare validator finds no violation of
Theorem~\ref{thm:subsquare-monotonicity}.  The multiples of four in the last
distribution also follow directly: an order-$m$ subsquare inside an
order-$2m$ square forces the three complementary blocks to be Latin
subsquares, so half-order subsquares occur in complementary quartets.

The ordered non-symmetric labels refer to the representatives in the cited
ANU file.  Parastrophy permutes the three bits, so labels such as
$\mathsf{FFT}$ are not themselves main-class invariants.  The FFF count, the
existence of every pattern, and the group-isotopy split are unaffected.

The scan tests every unordered pair of lines in all three views.  A separate
closure-family computation reproduces the $5/225$ split using
Proposition~\ref{prop:closure}.  The frozen ANU input is an exact byte match to
the file served by the cited URL on August 9, 2026.

\begin{corollary}\label{cor:all-patterns}
For every $m\geq3$, Latin squares of order $2^m$ realize all eight patterns.
Every such order also contains a non-group-isotopic FFF Latin square.
\end{corollary}

\begin{proof}
Computational Theorem~\ref{cthm:order8} supplies one order-$8$ square of each
pattern.  Take direct products with the FFF group table of $C_2^{m-3}$ and use
Theorem~\ref{thm:pattern-product}.  For the second assertion, start from the
explicit loop in Theorem~\ref{thm:explicit-nongroup-loop}; non-group-isotopy
is preserved by Theorem~\ref{thm:group-product}.
\end{proof}

This is an existence result only.  It neither counts higher-order main
classes nor asserts that products of distinct order-$8$ classes remain
distinct.

\noindent\textbf{Historical reproducibility record.} This section describes the underlying research collection, not a promise that every named artifact is in this public export. See the dossier scope statement and the public reproducibility guide for availability.
\medskip

\section{Computational reproducibility}\label{sec:reproducibility}

The historical research collection separates proofs, frozen data, programs, commands, logs, and
SHA-256 manifests.  A companion reproducibility quickstart supplies the
complete command-to-output map.  The compact paper uses the following
packages.

\begin{description}
\item[Trade parity]
\nolinkurl{repro_runs/2026-08-10_paper_symbol_sign_audit/}.
The exact control checks $1028307$ row-, column-, and symbol-view trades on all
reduced Latin squares of orders $2$, $4$, and $6$, with zero failures.  The
proof, not this computation, establishes the general theorem.

\item[Small orders]
\nolinkurl{repro_runs/2026-04-26_small_order_witness_reruns/}.
The exact-cover generator freshly enumerates $1$, $4$, and $9408$ reduced
squares of orders $2$, $4$, and $6$.  Every table is checked Latin, reduced,
and unique.  The generated order-$6$ set agrees exactly with the official ANU
file.

\item[Order-$8$ census]
\nolinkurl{repro_runs/2026-04-26_order8_core/}.
The frozen $5140393$-byte input has SHA-256
\nolinkurl{e5980ff5bc4ac46cbbefd1f1333258d0905724c42a976c4afe61ea501a4ccf87}.
The theorem-relevant output
\nolinkurl{results/order8_recheck_fullscan_results.json} has SHA-256
\nolinkurl{10f2f43c8dcaa8ec93f83abf6b768efa47f909a9c8623d4bc648e6da097ccd39}.

\item[Group-isotopy closure]
\nolinkurl{repro_runs/2026-04-26_order8_closure_reconstruction/}.
This independently reconstructs the basis-family closure scan and reproduces
the $5/225$ split.

\item[Direct products]
\nolinkurl{repro_runs/2026-04-26_direct_product_fff/}.
The validator checks the product-permutation formula in all three views and
tests the group-isotopy equivalence on group and non-group examples.

\item[Subsquare and congruence obstructions]
\nolinkurl{repro_runs/2026-08-09_fff_subsquare_congruence_obstructions/}.
The package scans all subsquares of the $230$ frozen FFF representatives,
exhausts the binary twists used as controls, and validates the explicit
twisted order-$8$ loop.

\item[Binary line kernel and half-order subsquares]
\nolinkurl{repro_runs/2026-08-11_fff_binary_homotopy_subsquares/}.
The package compares the subsquares constructed from every nonuniversal
binary line dependency with an independent direct enumeration on $9778$
tables.  All predicted and observed sets agree exactly, and the semantic
verifier reproduces the frozen order-$8$ distribution with zero errors.

\item[Prime-field line quotients]
\nolinkurl{repro_runs/2026-08-11_fff_prime_field_line_quotients/}.
The audit checks five prime fields on $9778$ tables, verifies all $1160$
nonuniversal quotient classes, and independently reconstructs $48890$
modular ranks with zero errors.  The theorem itself is proved separately.

\item[External-input release audit]
\nolinkurl{repro_runs/2026-08-09_paper_release_audit/}.
It downloads the cited ANU and Wanless inputs afresh and verifies exact byte
matches to the frozen files.  As with any computational theorem based on an
external census, completeness of the published enumeration remains a cited
scientific dependency.
\end{description}

Every package contains a README, exact commands, environment record, input
and output manifests, and a package-level SHA-256 manifest.  The command
\begin{center}
\small\texttt{python3 }\nolinkurl{repro_runs/2026-08-10_paper_symbol_sign_audit/scripts/audit_manuscript_computational_claims.py}
\end{center}
cross-checks the computational statements in the long companion manuscript;
the three computational theorems retained here are a subset of that audit.

Solver binaries, virtual environments and non-research source material are not included in the public export.  No timeout, partial solver model,
or order-$10$ diagnostic is used as evidence in this compact paper.

\section{Discussion and open problems}

Kobayashi and Nakamura studied the stricter condition that every pair of
factors in a $1$-factorization of $K_{n,n}$ forms only $4$-cycles; such a
factorization exists precisely in power-of-two orders~\cite{KobayashiNakamura1994}.
Our F condition permits arbitrary even induced cycle lengths, but requires the
condition simultaneously in all three coordinate views.  The neighbouring
classification, the odd-order obstruction, the order-$6$ computation, and the
power-of-two constructions motivate the following conjecture.

\begin{conjecture}[Power-of-two FFF conjecture]\label{conj:power-two}
For $n>1$, an FFF Latin square of order $n$ exists if and only if $n$ is a
power of two.
\end{conjecture}

The positive direction follows from Corollary~\ref{cor:group-fff}.  The
negative direction is proved for odd $n$ by Proposition~\ref{prop:odd-order}
and computed for $n=6$ by Computational Theorem~\ref{cthm:small-orders}.
The accompanying complete report excludes order $10$ by exhaustive
cycle-palette master searches.  The first undecided even non-power-of-two
order is therefore $12$.

Corollary~\ref{cor:2p-simple} provides a structural reduction at order $10$:
every order-$10$ FFF candidate would have to be congruence-simple
in every isotope.  Hence no central extension, binary extension, or other
congruence-imprimitive construction can settle the positive direction.  This
compact-paper theorem alone does not exclude simple quasigroups; the companion
report's exhaustive computation does.

Three concrete questions remain:
\begin{enumerate}
\item Decide the conjecture at order $12$ or prove a uniform obstruction for
even non-powers of two.
\item Find an order-independent obstruction coupling all three link
factorizations more strongly than their common intercalate count.
\item Determine when direct products preserve distinctions between isotopy or
main classes.  No such injectivity is asserted here.
\end{enumerate}

These questions concern one local odd-cycle obstruction.  They must not be
conflated with the global Alon--Tarsi conjecture, whose resolution would
require substantially more than existence or absence of this trade
primitive~\cite{AlonTarsi1992}.


\footnotesize
\bibliographystyle{alpha}
\bibliography{references}

\end{document}
